\documentclass[12pt]{article}
\usepackage[utf8]{inputenc}
\usepackage[T1]{fontenc}
\usepackage[margin=1in]{geometry}
\usepackage{setspace}
\usepackage{xcolor}
\usepackage[colorlinks=true, urlcolor=blue!60!black, linkcolor=blue!60!black]{hyperref}
\setlength{\parindent}{0em}
\setlength{\parskip}{0.45em}
\renewcommand{\baselinestretch}{1.0}
\begin{document}
\begin{center}
{\Large \textit{Virginia State Board of Pharmacy v. Virginia Citizens Consumer Council}, 425 U.S. 748 (1976)\par}
\vspace{0.3em}
{\normalsize Case Brief --- One-Page Edition\par}
\end{center}

\textbf{Memory Jogger.} Commercial speech gets First Amendment protection (\href{https://supreme.justia.com/cases/federal/us/425/748/}{425 U.S.\ 748}).

\textbf{Facts.} Virginia made it ``unprofessional conduct'' for pharmacists to advertise prescription-drug prices. Consumers sued, showing that prices for the same drugs varied widely within a single locality.

\textbf{Procedural History.} A three-judge district court struck the ban; the Board appealed directly.

\textbf{Votes.} Affirmed, 7--1 (Blackmun for the Court; Burger and Stewart concurring; Rehnquist dissenting).

\textbf{Holding.} Consumers have a reciprocal right to receive truthful advertising from willing speakers. Purely commercial speech is protected, and a State may not suppress truthful information about lawful activity out of fear of its effects.

\textbf{Reasoning.} Speech that does ``no more than propose a commercial transaction'' serves consumers, especially ``the poor, the sick, and particularly the aged,'' and society's interest in informed economic decisions. Virginia's ``highly paternalistic approach'' kept the public ignorant. False or misleading ads remain regulable. Rehnquist accused the Court of reviving \textit{Lochner}.

\textbf{Examples.} \emph{Same-Side:} lawyers' flat-fee ads (\textit{Bates}). \emph{Opposite-Side:} in-person solicitation (\textit{Ohralik}). \emph{Fence-Sitter:} limits on sportsbook bonus ads.

\textbf{Critique.} Critics see constitutionalized laissez-faire; Rehnquist foresaw direct-to-consumer drug ads.

\textbf{Key Quotation.} ``So long as we preserve a predominantly free enterprise economy, the allocation of our resources in large measure will be made through numerous private economic decisions. It is a matter of public interest that those decisions, in the aggregate, be intelligent and well informed. To this end, the free flow of commercial information is indispensable.'' --- Justice Blackmun, 425 U.S.\ at 765.

\end{document}
